% ============================================================
% Paper 95 (EN): Uniformization in Higher Dimensions:
%                Yau's Conjecture and Complex Geometry
% Author: Michael Fuhrmann
% Project: Specialist Mathematics Project
% Build: 169
% Date: 2026-03-12
% ============================================================
\documentclass[12pt,a4paper]{amsart}

\usepackage{amsmath,amssymb,amsthm,mathtools}
\usepackage{enumitem}
\usepackage{booktabs}
\usepackage{array}
\usepackage[hidelinks]{hyperref}
\usepackage[utf8]{inputenc}
\usepackage[T1]{fontenc}
\usepackage{geometry}
\geometry{margin=2.5cm}

% ---- Theorem Environments -----------------------------------
\newtheorem{theorem}{Theorem}[section]
\newtheorem{lemma}[theorem]{Lemma}
\newtheorem{corollary}[theorem]{Corollary}
\newtheorem{proposition}[theorem]{Proposition}
\theoremstyle{definition}
\newtheorem{definition}[theorem]{Definition}
\newtheorem{example}[theorem]{Example}
\theoremstyle{remark}
\newtheorem{remark}[theorem]{Remark}
\newtheorem{conjecture}[theorem]{Conjecture}

% ---- Custom Operators ---------------------------------------
\DeclareMathOperator{\Ric}{Ric}
\DeclareMathOperator{\Bsc}{Bsc}
\DeclareMathOperator{\Hol}{Hol}
\DeclareMathOperator{\vol}{vol}
\DeclareMathOperator{\diam}{diam}
\DeclareMathOperator{\rank}{rank}
\newcommand{\CC}{\mathbb{C}}
\newcommand{\RR}{\mathbb{R}}
\newcommand{\ZZ}{\mathbb{Z}}
\newcommand{\PP}{\mathbb{P}}
\newcommand{\NN}{\mathbb{N}}
\newcommand{\dd}{\partial}

% ---- Abstract -----------------------------------------------
\begin{document}
\begin{abstract}
The classical uniformization theorem classifies all Riemann surfaces up to
biholomorphism.  In higher dimensions, analogous classification results for
Kähler manifolds with positive curvature are provided by the theorem of
Mori (compact case, 1979) and Siu--Yau (compact Kähler, 1980): a compact
Kähler manifold with positive bisectional curvature is biholomorphic to
$\PP^n$.  Yau's conjecture---that a \emph{complete} (non-compact) Kähler
manifold with positive bisectional curvature is biholomorphic to $\CC^n$---
remains open in full generality.  We survey Cao's Kähler--Ricci flow
approach, pseudoconvexity, Stein manifolds, and related open problems.
\end{abstract}

\title{Uniformization in Higher Dimensions:\\
\large Yau's Conjecture and Complex Geometry}

\author{Michael Fuhrmann}
\address{Specialist Mathematics Project, Build 169}
\email{specialist-maths@project.local}
\date{2026-03-12}

\maketitle
\tableofcontents

% =============================================================
\section{Introduction}
% =============================================================

The \emph{uniformization theorem} for Riemann surfaces, proved by Poincaré
and Koebe, is one of the cornerstones of complex analysis:

\begin{theorem}[Poincaré--Koebe Uniformization, 1907]
Every simply-connected Riemann surface is biholomorphic to exactly one of:
\begin{enumerate}[label=(\roman*)]
  \item the Riemann sphere $\hat{\CC} \cong S^2$,
  \item the complex plane $\CC$,
  \item the upper half-plane $\mathbb{H} = \{z \in \CC : \Im z > 0\}$.
\end{enumerate}
\end{theorem}

This is a profound result: it says that the complex structure is determined
by topology (genus 0) together with completeness conditions.  In higher
dimensions, the situation is vastly more complex.

The natural higher-dimensional analogue asks: when is a complete Kähler
manifold with positive curvature biholomorphic to $\CC^n$ or $\PP^n$?
Yau formulated this as a conjecture in the 1970s, and the compact case was
fully resolved.  The non-compact case remains one of the central open
problems in complex geometry.

% =============================================================
\section{Kähler Manifolds and Bisectional Curvature}
% =============================================================

\begin{definition}[Kähler Manifold]
A complex manifold $(M, J)$ of complex dimension $n$ is \emph{Kähler} if
it admits a Hermitian metric $g$ whose associated $(1,1)$-form
$\omega(X,Y) = g(JX, Y)$ is closed: $d\omega = 0$.  Equivalently, the
Levi-Civita connection of $g$ is $J$-invariant: $\nabla J = 0$.
\end{definition}

\begin{example}
\begin{enumerate}[label=(\roman*)]
  \item $\CC^n$ with the Euclidean metric (flat Kähler).
  \item $\PP^n$ with the Fubini--Study metric (positive sectional curvature).
  \item Complex tori $\CC^n/\Lambda$ (flat Kähler, $\Lambda$ a lattice).
  \item Smooth projective varieties (Kodaira embedding + Fubini--Study restriction).
\end{enumerate}
\end{example}

\begin{definition}[Bisectional Curvature \cite{GoldbergKobayashi1967}]
Let $(M, g, J)$ be a Kähler manifold with curvature tensor $R$.  For unit
tangent vectors $X, Y \in T_p M$, the \emph{holomorphic bisectional curvature}
is
\[
  \Bsc(X,Y) = R(X, JX, Y, JY).
\]
The manifold has \emph{positive bisectional curvature} if
$\Bsc(X,Y) > 0$ for all non-zero $X, Y \in T_p M$ and all $p \in M$.
\end{definition}

\begin{remark}
Positive bisectional curvature implies positive sectional curvature, which
in turn implies positive Ricci curvature.  The implications are strict:
$\CC^n$ has zero bisectional curvature; the Fubini--Study metric on $\PP^n$
has constant positive bisectional curvature.
\end{remark}

% =============================================================
\section{The Classical Uniformization Theorem}
% =============================================================

We revisit the proof structure to identify what may generalise.

\begin{proof}[Sketch of Uniformization]
The proof proceeds in two steps:

\medskip
\noindent\textbf{Step 1 (Existence of harmonic maps).}
For a Riemannian surface $(\Sigma, g)$, construct a harmonic function
solving a Dirichlet problem to produce a conformal map to one of the
three models.

\medskip
\noindent\textbf{Step 2 (Curvature characterisation).}
Equivalently, the three models correspond to constant curvature $+1$
(sphere), $0$ (plane), $-1$ (half-plane). By the Gauss--Bonnet theorem,
the Euler characteristic determines which model applies:
$\chi > 0 \Rightarrow S^2$, $\chi = 0 \Rightarrow \CC$ (or torus),
$\chi < 0 \Rightarrow \mathbb{H}$.
\end{proof}

The higher-dimensional analogue of curvature $+1$ for Kähler manifolds is
\emph{positive bisectional curvature}, which motivates Yau's conjecture.

% =============================================================
\section{Mori's Bend-and-Break Theorem (Compact Case)}
% =============================================================

\begin{theorem}[Mori, 1979 \cite{Mori1979}]
\label{thm:mori}
Every smooth compact Kähler manifold $(M, g)$ with positive bisectional
curvature is biholomorphic to $\PP^n$ (for some $n = \dim_\CC M$).
\end{theorem}

Mori's proof is algebraic and uses the following key tool:

\begin{theorem}[Bend-and-Break \cite{Mori1979}]
Let $M$ be a smooth projective variety of dimension $n$ over an algebraically
closed field.  If $M$ contains a smooth curve $C$ with
$(-K_M \cdot C) > 0$ (i.e., $M$ has positive anti-canonical class in the
direction of $C$), then $M$ contains a rational curve.
\end{theorem}

\begin{proof}[Proof sketch of Mori's Theorem~\ref{thm:mori}]
\textbf{Step 1.}  By the Kodaira embedding theorem, a compact Kähler manifold
with positive bisectional curvature is projective.

\textbf{Step 2.}  Positive bisectional curvature implies
$c_1(M) > 0$ (by the Chern--Weil formula and positivity).  Thus
$-K_M$ is ample.

\textbf{Step 3.}  Using characteristic-$p$ techniques (Frobenius), Mori
deforms a curve in $M$ (bend-and-break) to produce rational curves.
These cover $M$: every point lies on a rational curve.

\textbf{Step 4.}  By a theorem of Kobayashi--Ochiai, if $M$ has index
$i(M) = n+1$ (i.e., $-K_M = (n+1)H$ for some ample $H$), then
$M \cong \PP^n$. $\square$
\end{proof}

% =============================================================
\section{Siu--Yau's Theorem: The Differential-Geometric Proof}
% =============================================================

\begin{theorem}[Siu--Yau, 1980 \cite{SiuYau1980}]
\label{thm:siuyau}
Every compact Kähler manifold with positive bisectional curvature is
biholomorphic to $\PP^n$.
\end{theorem}

Siu--Yau's proof is differential-geometric and predates Mori's algebraic
proof (they were proved around the same time, 1979--1980).

\begin{proof}[Proof sketch]
\textbf{Step 1 (Frankel Conjecture).}  The result was conjectured by Frankel
(1961) as: "a compact Kähler manifold with positive bisectional curvature is
homeomorphic to $\PP^n$."

\textbf{Step 2 (Existence of harmonic map).}  Using the theory of harmonic
maps, Siu--Yau show the existence of a holomorphic map
$f \colon M \to \PP^n$ of degree 1, using Bochner--Weitzenböck formulas.
The positivity of bisectional curvature ensures rigidity.

\textbf{Step 3 (Biholomorphism).}  The degree-1 holomorphic map $f$, combined
with the classification of holomorphic maps $\PP^n \to \PP^n$, implies that
$f$ is a biholomorphism. $\square$
\end{proof}

\begin{remark}
The hypothesis of \emph{positive} bisectional curvature is essential.
Non-negative bisectional curvature leads to a weaker conclusion (Mok's
theorem, 1988: biholomorphic to a compact Hermitian symmetric space).
\end{remark}

% =============================================================
\section{Yau's Conjecture: The Non-Compact Case}
% =============================================================

\begin{conjecture}[Yau, 1970s \cite{Yau1977}, Open in Full Generality]
\label{conj:yau}
Let $(M, g)$ be a complete Kähler manifold with positive bisectional
curvature.  Then $M$ is biholomorphic to $\CC^n$.
\end{conjecture}

\begin{remark}
Here "complete" means geodesically complete (every geodesic extends to
$\RR$).  This is the natural completeness condition; without it,
open balls in $\CC^n$ with flat metric would be trivial counterexamples.
\end{remark}

The conjecture is known in special cases:

\begin{theorem}[Partial Results on Conjecture~\ref{conj:yau}]
\begin{enumerate}[label=(\roman*)]
  \item \textbf{Complex dimension 2 (Mok, 1984 \cite{Mok1984})}:
        A complete Kähler surface with positive bisectional curvature and
        finite volume is biholomorphic to $\CC^2$.
  \item \textbf{Bounded curvature (Shi, 1989 \cite{Shi1989})}:
        If additionally the curvature is bounded ($|R| \le C$), then
        $M \cong \CC^n$.
  \item \textbf{Maximum volume growth (Chen--Tang--Zhu, 2004 \cite{ChenTangZhu2004})}:
        If $M$ has maximal volume growth ($\vol(B(p,r)) \ge c r^{2n}$), then
        $M \cong \CC^n$.
  \item \textbf{Maximum volume growth — full proof (Gang Liu, 2016 \cite{GangLiu2016})}:
        Gang Liu gave a complete, self-contained proof that every complete
        Kähler manifold with non-negative bisectional curvature and
        \emph{maximal volume growth} is biholomorphic to $\CC^n$.  This
        removes the bounded-curvature hypothesis of Shi and simplifies
        the argument of Chen--Tang--Zhu.
\end{enumerate}
\end{theorem}

\begin{remark}
Liu's 2016 theorem \cite{GangLiu2016} fully resolves Yau's conjecture
in the \emph{maximal volume growth} case without any curvature-boundedness
assumption.  The remaining open case is when the volume growth is
\emph{sub-maximal} (i.e.\ $\vol(B(p,r)) = o(r^{2n})$ as $r \to \infty$)
combined with positive (not just non-negative) bisectional curvature.
\end{remark}

% =============================================================
\section{The Kähler--Ricci Flow Approach}
% =============================================================

Hamilton's Ricci flow \cite{Hamilton1982} deforms a Riemannian metric by
$\frac{\partial g}{\partial t} = -2\Ric(g)$.  For Kähler manifolds,
Cao \cite{Cao1985} showed the flow preserves the Kähler condition.

\begin{theorem}[Cao, 1985 \cite{Cao1985}]
The Kähler--Ricci flow
\[
  \frac{\partial g_{i\bar{j}}}{\partial t} = -R_{i\bar{j}}
\]
has a unique solution for all $t \in [0,\infty)$ starting from any compact
Kähler metric of positive bisectional curvature.  The flow converges (after
normalisation) to the Fubini--Study metric on $\PP^n$.
\end{theorem}

\begin{theorem}[Hamilton, 1988 \cite{Hamilton1988}]
In complex dimension 1 (Riemann surfaces), the normalised Kähler--Ricci flow
converges to a metric of constant curvature.  This gives an analytical proof
of the uniformization theorem for compact Riemann surfaces.
\end{theorem}

For the non-compact case:

\begin{theorem}[Shi, 1989 \cite{Shi1989}]
The Ricci flow on a complete non-compact Riemannian manifold with bounded
curvature has a unique solution for short time.  For complete Kähler manifolds
with bounded positive bisectional curvature, the flow preserves the Kähler
structure and remains complete.
\end{theorem}

The strategy for Yau's conjecture via Kähler--Ricci flow:
\begin{enumerate}[label=(\roman*)]
  \item Run Kähler--Ricci flow from $(M,g)$ with positive bisectional curvature.
  \item Show the flow exists for all time.
  \item Show the flow converges to the flat metric on $\CC^n$.
  \item Conclude $M \cong \CC^n$.
\end{enumerate}
Steps (i) and (ii) are known under bounded curvature assumptions;
steps (iii) and (iv) require additional work.

% =============================================================
\section{Stein Manifolds and Pseudoconvexity}
% =============================================================

Yau's conjecture is also related to the theory of Stein manifolds, which are
the natural "affine" objects in complex geometry.

\begin{definition}[Stein Manifold]
A complex manifold $M$ is \emph{Stein} if it satisfies:
\begin{enumerate}[label=(\roman*)]
  \item \textbf{Holomorphic convexity}: for every compact $K \subseteq M$,
        the holomorphic hull $\hat K_\OO = \{p \in M : |f(p)| \le \sup_K |f|
        \text{ for all holomorphic } f\}$ is compact.
  \item \textbf{Holomorphic separability}: for $p \neq q$ there exists $f$
        holomorphic with $f(p) \neq f(q)$.
  \item \textbf{Holomorphic generation of cotangent bundle}: at each point,
        $\{df_p : f \text{ holomorphic on } M\}$ spans $T^*_p M$.
\end{enumerate}
\end{definition}

\begin{theorem}[Grauert, 1958 \cite{Grauert1958}]
A complex manifold $M$ is Stein if and only if it admits a strictly
plurisubharmonic exhaustion function $\varphi \colon M \to \RR$ (a smooth
function with $i\dd\bar\dd\varphi > 0$ in the sense of $(1,1)$-forms,
with compact sublevel sets).
\end{theorem}

\begin{theorem}[Greene--Wu, 1978 \cite{GreeneWu1978}]
A complete Kähler manifold with positive bisectional curvature is Stein.
\end{theorem}

\begin{corollary}
Yau's conjecture would imply that the only Stein Kähler manifold with
positive bisectional curvature is $\CC^n$.
\end{corollary}

% =============================================================
\section{The Hyperconvexity Conjecture}
% =============================================================

\begin{definition}[Hyperconvex Domain]
A domain $\Omega \subseteq \CC^n$ is \emph{hyperconvex} if it admits a
bounded plurisubharmonic exhaustion function $u \colon \Omega \to [-\infty, 0)$
(i.e., $u$ is plurisubharmonic, bounded above, and $\{u \le c\} \Subset \Omega$
for all $c < 0$).
\end{definition}

\begin{conjecture}[Hyperconvexity Conjecture, Open]
\label{conj:hyperconvex}
Every pseudoconvex domain in $\CC^n$ with $C^1$ boundary is hyperconvex.
\end{conjecture}

\begin{remark}
Hyperconvex domains have especially nice properties: the Dirichlet problem
for the complex Monge--Ampère equation
$(i\partial\bar\partial u)^n = f \cdot \omega^n$
is solvable, and the Bergman kernel behaves well.
\end{remark}

\begin{theorem}[Diederich--Fornæss, 1977 \cite{DiederichFornaess1977}]
Every bounded pseudoconvex domain with $C^2$ boundary in $\CC^n$ is hyperconvex.
\end{theorem}

The conjecture for $C^1$ domains remains open.  Counterexamples exist for
merely Lipschitz domains (Worm domains).

% =============================================================
\section{Mok's Theorem on Non-Negative Bisectional Curvature}
% =============================================================

\begin{theorem}[Mok, 1988 \cite{Mok1988}]
\label{thm:mok}
Let $(M, g)$ be a compact irreducible Kähler manifold with non-negative
bisectional curvature.  Then $M$ is biholomorphic to one of the following
\emph{compact irreducible Hermitian symmetric spaces}:
\begin{enumerate}[label=(\roman*)]
  \item $\PP^n$ (complex projective space),
  \item $\mathrm{Gr}(k,n)$ (complex Grassmannian),
  \item $\mathrm{SO}(n)/U(n/2)$ (type DIII),
  \item $\mathrm{Sp}(n)/U(n)$ (type CI),
  \item $E_6/\mathrm{Spin}(10)\cdot U(1)$,
  \item $E_7/E_6 \cdot U(1)$.
\end{enumerate}
\end{theorem}

\begin{remark}
Mok's theorem is the full classification under the weaker hypothesis of
non-negative (rather than strictly positive) bisectional curvature.
It confirms that positive curvature is far more rigid than non-negative.
\end{remark}

% =============================================================
\section{The Riemann Surface Case Revisited}
% =============================================================

For Riemann surfaces, the uniformization theorem classifies all surfaces.
The higher-dimensional analogue has distinct compact and non-compact versions:

\begin{center}
\begin{tabular}{p{3cm}p{3.5cm}p{3.5cm}p{3.5cm}}
\toprule
 & Riemann surface & Compact Kähler & Complete Kähler \\
\midrule
Positive curvature & $S^2 \cong \PP^1$ & $M \cong \PP^n$ (Mori--Siu--Yau, proved) & $M \cong \CC^n$? (Yau, open) \\[4pt]
Zero curvature & $\CC$ or $\TT^2$ & Kähler flat (CY) & Flat ($M \cong \CC^n/\Gamma$) \\[4pt]
Negative curvature & $\mathbb{H}$ & $M$ = general type & Cartan--Hadamard \\
\bottomrule
\end{tabular}
\end{center}

% =============================================================
\section{Open Problems}
% =============================================================

\begin{conjecture}[Yau's Conjecture on Positive Bisectional Curvature,
= Conjecture~\ref{conj:yau}, Open in Full Generality]
Every complete Kähler manifold with positive bisectional curvature is
biholomorphic to $\CC^n$.
\end{conjecture}

\begin{conjecture}[Yau's Uniformization Conjecture for Positive Ricci, Open]
Every complete Kähler manifold with positive Ricci curvature, bounded
curvature, and maximal volume growth is biholomorphic to $\CC^n$.
\end{conjecture}

\begin{conjecture}[Frankel Conjecture (Non-Compact Version), Open]
Every complete simply-connected Kähler manifold with positive sectional
curvature is biholomorphic to $\CC^n$.
\end{conjecture}

\begin{conjecture}[Hyperconvexity for $C^1$ Domains,
= Conjecture~\ref{conj:hyperconvex}, Open]
Every pseudoconvex domain in $\CC^n$ with $C^1$ boundary is hyperconvex.
\end{conjecture}

\begin{conjecture}[Grauert Tube Conjecture, Open]
The Grauert tube of a real-analytic Riemannian manifold with positive
sectional curvature is biholomorphic to a ball in $\CC^n$.
\end{conjecture}

% =============================================================
\section{Conclusion}
% =============================================================

The uniformization theorem for Riemann surfaces has inspired decades of
effort to understand complex manifolds in higher dimensions.  The compact
case of Yau's conjecture was resolved by Mori and Siu--Yau: compact Kähler
manifolds with positive bisectional curvature are $\PP^n$.

The non-compact case---Yau's conjecture proper---remains open in full
generality.  Partial results under curvature bounds (Shi, 1989) and
maximum volume growth (Chen--Tang--Zhu, 2004; Gang Liu, 2016) give evidence
for the conjecture.  The Kähler--Ricci flow approach (Cao, 1985) provides a
potentially powerful tool, but the flow's long-time behaviour in the
non-compact case without bounded curvature is not fully understood.

The connections to pseudoconvexity, Stein manifolds, and the Monge--Ampère
equation make this an active intersection of complex analysis, differential
geometry, and partial differential equations.

\begin{thebibliography}{99}

\bibitem{Yau1977}
S.-T. Yau,
\textit{Calabi's conjecture and some new results in algebraic geometry},
Proc. Nat. Acad. Sci. USA \textbf{74} (1977), 1798--1799.

\bibitem{Mori1979}
S.~Mori,
\textit{Projective manifolds with ample tangent bundles},
Ann. of Math. (2) \textbf{110} (1979), 593--606.

\bibitem{SiuYau1980}
Y.-T. Siu and S.-T. Yau,
\textit{Compact Kähler manifolds of positive bisectional curvature},
Invent. Math. \textbf{59} (1980), 189--204.

\bibitem{Mok1984}
N.~Mok,
\textit{An embedding theorem of complete Kähler manifolds of positive bisectional curvature onto affine algebraic varieties},
Bull. Soc. Math. France \textbf{112} (1984), 197--258.

\bibitem{Mok1988}
N.~Mok,
\textit{The uniformization theorem for compact Kähler manifolds of nonnegative holomorphic bisectional curvature},
J. Differential Geom. \textbf{27} (1988), 179--214.

\bibitem{Cao1985}
H.-D. Cao,
\textit{Deformation of Kähler metrics to Kähler--Einstein metrics on compact
Kähler manifolds},
Invent. Math. \textbf{81} (1985), 359--372.

\bibitem{Hamilton1982}
R.~S. Hamilton,
\textit{Three-manifolds with positive Ricci curvature},
J. Differential Geom. \textbf{17} (1982), 255--306.

\bibitem{Hamilton1988}
R.~S. Hamilton,
\textit{The Ricci flow on surfaces},
Contemp. Math. \textbf{71} (1988), 237--262.

\bibitem{Shi1989}
W.-X. Shi,
\textit{Ricci deformation of the metric on complete noncompact Riemannian
manifolds},
J. Differential Geom. \textbf{30} (1989), 223--301.

\bibitem{ChenTangZhu2004}
B.-L. Chen, S.-H. Tang, and X.-P. Zhu,
\textit{A uniformization theorem for complete non-compact Kähler surfaces
with positive bisectional curvature},
J. Differential Geom. \textbf{67} (2004), 519--570.

\bibitem{GoldbergKobayashi1967}
S.~I. Goldberg and S.~Kobayashi,
\textit{Holomorphic bisectional curvature},
J. Differential Geom. \textbf{1} (1967), 225--233.

\bibitem{Grauert1958}
H.~Grauert,
\textit{On Levi's problem and the imbedding of real-analytic manifolds},
Ann. of Math. (2) \textbf{68} (1958), 460--472.

\bibitem{GreeneWu1978}
R.~E. Greene and H.~Wu,
\textit{$C^\infty$ approximations of convex, subharmonic, and plurisubharmonic
functions},
Ann. Sci. Éc. Norm. Supér. (4) \textbf{12} (1979), 47--84.

\bibitem{GangLiu2016}
G.~Liu,
\textit{Compactness and finiteness theorems for Ricci flow and K\"ahler--Ricci flow with bounded scalar curvature},
and
\textit{Three-circle theorem and dimension estimate for holomorphic functions on K\"ahler manifolds},
Duke Math.\ J.\ \textbf{165} (2016), no.~15, 2899--2919.

\bibitem{DiederichFornaess1977}
K.~Diederich and J.~E. Fornæss,
\textit{Pseudoconvex domains: bounded strictly plurisubharmonic exhaustion
functions},
Invent. Math. \textbf{39} (1977), 129--141.

\end{thebibliography}

\end{document}
